\documentclass[11pt]{article}
\usepackage[margin=1in]{geometry}
\usepackage{amsmath,amssymb,hyperref}
\newcommand{\F}{\mathbb F}
\title{One weighted convolution for the normalized section product}
\author{Project research note}
\date{18 November 2026}
\begin{document}
\maketitle

The section product in \texttt{../implicit\_resultant/resultant.tex}
admits a joint formula on the odd prime-order family. It gives the
entire normalized polynomial in the target coordinate, including its
leading scalar. Two additional integer moments in the convolution
give the full additive gradient at every zero, even when distinct
sign choices give the same zero tuple. A dense implementation costs
$2^{b+o(b)}$ with preprocessing charged. It can share that cost across
all targets of a seed. The previous single-target and end-to-end bounds
remain better. We give a restricted lower bound for compressing the
new unit product, using the actual dyadic masks.

\section{Removing all intermediate section scales}
Let $E:y^2=x^3+ax+c$ be smooth over $\F_p$, $p>3$. Work in a cyclic
group $G$ of odd order $r$ with $\gcd(r,p-1)=1$. For the asymptotic
statements take $G=E(\F_p)$ of odd prime order. Large prime-order
curves satisfy the coprimality since $r>(p-1)/2$ by Hasse and $p-1$
is even. Use $k(P)=(x(P),1)$ for $P\ne O$ and $k(O)=(1,0)$.
Write $\ell_P(Z)=Z_0 k_1(P)-Z_1 k_0(P)$, so
$\ell_O(Z)=-Z_1$. Keeping this sign is necessary.
The homogeneous chord form $K$ of the preceding note factors as
\[
 K(k(P),k(Q),Z)=L(P,Q)\,\ell_{P+Q}(Z)\ell_{P-Q}(Z).
\]
The multiplier is $1$ if either input is $O$, $(x(P)-x(Q))^2$
if the finite abscissas differ, and $4y(P)^2$ if they coincide.
The coincident case follows by taking the linear coefficient of the
affine section and accounting for $\ell_O=-1$.

There is an even function $\gamma:G\to\F_p^\times$, $\gamma(O)=1$,
for which
\begin{equation}
 L(P,Q)=\frac{\gamma(P)^2\gamma(Q)^2}{\gamma(P+Q)\gamma(P-Q)}.
 \label{scale}
\end{equation}
Here is a way to obtain the absolute scale. Let $\Delta_r$ be a
homogeneous lift of multiplication by $r$ on the abscissa line, of
degree $r^2$, with $\Delta_r(1,0)=(1,0)$. These are the usual
coprime projective multiplication polynomials, normalized at infinity.
See the construction in
\href{https://ocw.mit.edu/courses/18-783-elliptic-curves-fall-2025/mit18_783_f25_lec_full.pdf}
{Sutherland's elliptic curve notes, Theorem 5.25}.
For $P\in G$ write $\Delta_r(k(P))=d(P)k(O)$.
Applying the chord map and then applying $\Delta_r$ to its two
output roots gives the same homogeneous binary quadratic as applying
$\Delta_r$ to both inputs and then the chord map. To see equality of
scales, the induced map on an unordered root pair sends
$\ell_U\ell_V$ to $\ell_{\Delta_r(U)}\ell_{\Delta_r(V)}$ in vector
coordinates. It has degree $r^2$ in the quadratic coefficients
(equivalently use the homogeneous resultant with
$\det(Z,\Delta_r(W))$). Both composites are base-point-free forms
of bidegree $(2r^2,2r^2)$. They agree projectively by
$[r](P\pm Q)=[r]P\pm[r]Q$, hence differ by a global nonzero
constant. Evaluation at $(O,O)$ fixes it to one. Therefore
\[
 L(P,Q)^{r^2}d(P+Q)d(P-Q)=d(P)^2d(Q)^2.
\]
Set $\gamma(P)=d(P)^{(r^2)^{-1}\bmod(p-1)}$.
In particular $\gamma(2P)=\gamma(P)^4/(4y(P)^2)$ for finite $P$.
Doubling permutes $G$, so every $\gamma(P)$ is a square.

The finite implementation also computes these scales without expanding
division polynomials. For a generator $B$ let $v_0=v_1=1$ and
$v_{i+1}=v_i^2/(v_{i-1}L(iB,B))$ up to $i=r-1$.
Equation~\eqref{scale} gives
$\gamma(iB)=v_i\,\gamma(B)^{i^2}$ by induction. The endpoint
$\gamma(rB)=1$ fixes $\gamma(B)=v_r^{-(r^2)^{-1}\bmod(p-1)}$.
The table of group elements and scales costs $O(r\,\mathrm{poly}(b))$.

\section{The actual joint product}
For each nonempty paired slot $U_i$ choose one lift $P_u$ of each
abscissa. Fix the first sign to $+$ and allow both signs in every
other slot. Let $d(S)$ count these oriented tuples whose sum is
$S$. This is one ordinary integer convolution of the first
oriented mask and the remaining signed masks. Define
\[
 D=2^{m-1}\prod_i|U_i|,\qquad
 A=\prod_i\left(\prod_{u\in U_i}\gamma(P_u)\right)^{D/|U_i|}.
\]
Then the exact normalized product of the project's summation forms is
\begin{equation}
 N(Z,\mathbf0)=A\prod_{S\in G}
   \left(\frac{\ell_S(Z)}{\gamma(S)}\right)^{d(S)}.
 \label{joint}
\end{equation}
For one input the $\gamma$ factors cancel. When a new abscissa
is appended, the homogeneous resultant squares the existing
leading scalar and replaces each $\ell_S$ by
$L(S,P)\ell_{S+P}\ell_{S-P}$. Substitution of~\eqref{scale}
cancels the old $\gamma(S)$ factors. This proves the single-tuple
identity by induction. Multiplication over all abscissa tuples
gives~\eqref{joint}. With an empty slot $N=1$. No intermediate
section, half polynomial, or discrete-log oracle is needed.
The implementation obtains $U_i$ directly from the dyadic
inverse hash, intersecting canonical abscissas with its charged
curve table.

Collecting equal finite abscissas gives $h_x=d(S)+d(-S)$ at
$x=x(S)$. In affine coordinates write
\[
 C=A(-1)^{d(O)}\prod_{S\in G}\gamma(S)^{-d(S)},\qquad
 N(z,\mathbf0)=C\prod_x(z-x)^{h_x}.
\]
For an eligible target $t=x(T)$, set $h=h_t$, taking zero for a
missing entry, and define its absolute unit
\begin{equation}
 U_t=C\prod_{x\ne t}(t-x)^{h_x}.
 \label{unit}
\end{equation}
This is the coefficient of $(z-t)^h$ in the local expansion.
It is the norm at order zero and $N_z(t)$ at order one.
At higher order $N_z(t)=0$. There is no factorial in this
formula, even if $h\ge p$.

\section{First derivatives and many targets}
Encode an abscissa tuple injectively by the ordinary integer
$q=\sum_{i=1}^m u_i p^{i-1}$ using canonical $u_i\in[0,p)$.
At each group element keep the count, the sum of $q$, and the sum
of $q^2$. Under multiplication the update for the last component is
$M_2*N'+2M_1*M'_1+N*M'_2$, where each star is cyclic convolution,
and the first moment is $M_1*N'+N*M'_1$. This follows from adding
the codes of disjoint sets of slots. Fold the three entries at
$T,-T$ to get $(h,M_1,M_2)$.
For $h>0$, $hM_2-M_1^2$ is the sum of squared differences over
all unordered pairs of codes. Exact integer arithmetic therefore
distinguishes multiple abscissa tuples from a single tuple of
arbitrary sign multiplicity.

If the variance is positive, at least two different factors of
$N$ vanish and every first derivative is zero. Otherwise $M_1/h$
decodes the unique tuple $\mathbf u^*$. For its single normalized
form $S_*(z)$, compute $v_*=(S_*(z)/(z-t)^h)|_{z=t}$, which is
nonzero. With the additive marker convention
$S(\mathbf u;z)+s_0+\sum_i s_i u_i$ one obtains
\[
 N_{s_0}=U_t/v_*,\qquad N_{s_i}=u_i^*\,U_t/v_*.
\]
In particular all nonzero sections that set the cofactor scale
are included. This recovers the full gradient at a zero in one
joint convolution with two moment channels. At a nonzero norm,
$N_z$ follows by logarithmic differentiation of~\eqref{joint}.
Literal additive derivatives there still involve sums weighted by
$1/S(\mathbf u;t)$. The two-slot algorithm in the preceding
scalar note supplies them where requested. The relation query
can exit on a nonzero norm.

For many targets under one seed, put
$g_j(X)=\prod_{x:\operatorname{bit}_j(h_x)=1}(X-x)$. At a query
$t$, use $g_j(t)$ when nonzero and $g'_j(t)$ otherwise.
Each $g_j$ is squarefree, including in positive characteristic.
Raising these values to $2^j$, multiplying over $j$ and multiplying
by $C$ gives $U_t$. At a nonzero, the logarithmic derivative is
$\sum_j2^jg'_j(t)/g_j(t)$. A subproduct and remainder tree shares
all these evaluations across the batch. With $r$ table entries,
at most $r$ targets and integer convolution coefficients of
$O(mb+\log D)$ bits, the work is
$r\,2^{O(m)}\mathrm{poly}(b,m)$ including the single-tuple forms
at zeros. Memory is $r\,2^{O(m)}\mathrm{poly}(b,m)$. The same
bound covers every target of one seed. Already writing the group
table costs $\Omega(r)$ for this construction.

\section{A lower bound for shortcuts that separate the slots}
Consider a single bit of~\eqref{unit}. Let $\chi$ be the quadratic
character and write $(-1)^{k_T(S)}$ for the character of $t-x(S)$
at ordinary points, with value $+1$ at $\pm T$ and $\chi(-1)$ at
$O$. All $\gamma$ factors are squares, so
\[
 \chi(U_t)=(-1)^{\sum_{S\in G}d(S) k_T(S)} .
\]
This formula includes seeds with no zero and retains a bit of the
absolute derivative when $h=1$. With two singleton abscissa slots
the matrix of output bits, indexed by nonzero sign pairs, is
\[
 B_T(P,Q)=k_T(Q+P)+k_T(Q-P)\pmod2.
\]
Its rank over $\F_2$ is $\Omega(\sqrt p)$ at \emph{every} eligible
target on the large prime-order family.

For the proof consider a nonempty proper set $J$ of $s$ distinct
shifts. The rational function
$f_J(X)=\prod_{A\in J}(t-x(X+A))$ has an odd-order zero or pole
somewhere. In fact its divisor modulo two has support given by
the symmetric difference of $T-J$ and $-T-J$. Equality of these
sets would make $J$ invariant under translation by $2T$, a
generator. Thus the double cover $Y^2=f_J(X)$ is geometrically
connected. It has at most $2s$ branch points, so its smooth
projective model has genus at most $1+s$ by Riemann--Hurwitz.
Away from at most $3s$ exceptional points the number of rational
points in a fiber minus one equals the product of the shifted
signs. At the exceptions the discrepancy is at most two per
point. The Hasse--Weil bound for the cover and for $E$ gives
\[
 \left|\sum_{X\in G}(-1)^{\sum_{A\in J}k_T(X+A)}\right|
 \le (2s+4)\sqrt p+6s .
\]
For the curve bound see Weil's \emph{Sur les courbes alg\'ebriques
et les vari\'et\'es qui s'en d\'eduisent} (1948) and the statement in
\href{https://arxiv.org/abs/1409.6090}{van der Geer's survey}.
If the right side is less than $r$, such a parity of shifts cannot
be constant. A sum of $j$ distinct rows of $B_T$ uses $2j$
distinct shifts, is even as a function of $Q$, and is zero at $O$.
Hence any $j$ with $(4j+4)\sqrt p+12j<r$ gives independent
rows. This proves the rank bound.

For actual masks with fixed odd multiplier and varying shift, let
$n=(r-1)/2$ and let $\delta$ be the codimension of the span of
the abscissa masks in $\F_2^n$. The shift-span calculation in
\texttt{../slot\_arity/rate.tex} gives $\delta\le\tau$, and
$\delta=0$ if the omitted canonical tail meets all but at most one
residue class modulo $\tau$. For two slots the matrix over the
independent seed shifts is $M_L B_T M_R^\top$, where each $M$
has the actual masks as rows. Its rank is at least
$\operatorname{rank}B_T-\delta_L-\delta_R$.

The same bound holds for the full matrix across a nontrivial
partition of the growing collection of slots, leaving one free
slot on each side. Fix the first sign in the free slot on the
left. For every other slot except the two free ones, take a
nonzero difference of two adjacent shift masks. As a signed
group function it has the form $e_P$ or $e_P+e_Q$ over $\F_2$,
where $e_P=[P]+[-P]$ and the sign pairs in a sum are distinct.
Such a difference exists whenever $0<\tau<M$ and some canonical
abscissa occurs. On every nontrivial Fourier character of the
odd prime-order group its value is nonzero. Indeed $z+z^{-1}$
is nonzero for a nontrivial $r$th root, and equal values give
equal or inverse roots. Thus these extra convolutions preserve the
rank on the span of $e_P$, which in Fourier coordinates consists
of the independent nontrivial inversion pairs. Restricting the
two free masks costs at most their two codimensions as before.
The contractions are row and column operations on the full
seed matrix. Since the growing-arity widths are $p^{o(1)}$,
the resulting lower bound is still $\Omega(\sqrt p)$.
It prices exact sums of separated left/right products over
$\F_2$ for the unit bit on all inputs. It leaves open joint or
average-case computations and routines specialized to successful trials.

\section{Checks and end-to-end accounting}
Run \texttt{python -B research/section\_product/verify.py}.
On the order-19 curve all 63 dyadic subboxes plus three additional
inputs gave 594 exact comparisons of order, absolute unit and target
derivative against the literal product. The 114 zeros gave matching
full gradients, including six single tuples with repeated sign
completions and 29 cases with several zero tuples. The 2,800-tuple
fixture over $\F_{65521}$ gave $U_t=N_z=43446$ and
$N_{s_0}=16953$, with coordinate derivatives
$13292,47523,12119,9163$. The order, unit and target derivative
also matched an independent local remainder of the literal
2,800-factor polynomial at every one of the 32,709 target abscissas.
At all 15,934 zeros the full gradients matched a product-rule
reference carrying the five additive marker polynomials.
At $t=0$ the sign matrices attained rank 509 of 509 and 32709 of
32709 on the order-1019 and order-65419 curves. The smaller matrix
was checked directly and by a cyclic-polynomial gcd, with the
width-32 hash masks spanning all 509 abscissa coordinates. For
the larger curve the gcd alone gives the rank. Indeed the binary
circulant for $k_T$ has rank
$r-\deg\gcd(\sum_i k_T(iB)X^i,X^r-1)$, and its nontrivial Fourier
coordinates come in equal inversion pairs. Removing the trivial
component and halving gives the rank of $B_T$.

The construction aggregates the normalized section contributions
and shares their absolute derivatives at zeros across targets.
Its charged $2^{b+o(b)}$ preprocessing exceeds the
$2^{(1/2+o(1))b}$ explicit-halves query and the
$2^{(2/3+o(1))b}$ three-block query. The training budget at tuned
arity is $2^{b/m+o(b)}$ rows, so batching these queries does not
absorb the group-sized preprocessing. The best project estimate
continues to be
\[
 \log_2\mathbb E T_{\rm solve}\le
 (26.2933775156+o(1))\,b/(\log_2 b)^{1/4}
\]
under the project's finite-packet and contraction assumptions.
The remaining issue for this formula is to evaluate the
target-weighted product and the required target moments from the
independent seeds within that budget, without materializing the
group distribution.
\end{document}
